\documentclass[11pt]{amsart}
\usepackage[T1]{fontenc}
\usepackage{lmodern,amsmath,amssymb,amsthm,booktabs,array,xurl}
\usepackage[margin=1in]{geometry}
\usepackage[colorlinks=true,urlcolor=blue,linkcolor=blue,citecolor=blue]{hyperref}
\newcommand{\Z}{\mathbb Z}
\newcommand{\F}{\mathbb F}
\newcommand{\R}{\mathbb R}
\newcommand{\Q}{\mathbb Q}
\newcommand{\norm}[1]{\lVert#1\rVert}
\newcommand{\lean}[1]{{\small\nolinkurl{#1}}}
\newcommand{\src}[1]{\href{https://github.com/tomdif/RamanujanLean/blob/main/RamanujanTau/MultiQuintuple#1}{\lean{#1}}}
\newtheorem{theorem}{Theorem}
\newtheorem{proposition}[theorem]{Proposition}
\newtheorem{lemma}[theorem]{Lemma}
\theoremstyle{definition}
\newtheorem{definition}[theorem]{Definition}
\newtheorem{remark}[theorem]{Remark}
\title[Vanishing in triple quintuple products]{Persistent coefficient vanishing in admissible products of three quintuple products}
\author{Thomas DiFiore}
\date{Review draft, November 11, 2026}
\begin{document}
\begin{abstract}
For a prime modulus and three distinct canonical indices whose squared sum
is divisible by that modulus, we give a necessary and sufficient arithmetic
criterion for an entire coefficient progression of a product of three
quintuple products to vanish. The criterion is a short vector in a projective
dual lattice, together with an explicit residue equation. The equivalence
is proved in Lean. We explain the analytic and lattice steps, the finite
certificates below the geometric threshold, and a linear scan of projective
scalars that implements the criterion. We also give a separate, completely
classified four-factor example at modulus 97. This is a manuscript for
independent mathematical and priority review; it does not assert a general
four-factor converse.
\end{abstract}
\maketitle

\section{Statement and scope}
Write $(a;q)_\infty=\prod_{n\ge0}(1-aq^n)$ and fix the convention
\[
 Q(z,q)=(z,q/z,q;q)_\infty(qz^2,q/z^2;q^2)_\infty.
\]
Thus $Q(q^i,q^p)$ is an integral formal power series when $0<2i<p$.
Put
\[
 F_{p,v}(q)=Q(q^i,q^p)Q(q^j,q^p)Q(q^k,q^p)
       =\sum_{n\ge0}a_{p,v}(n)q^n,\qquad v=(i,j,k).
\]
\begin{definition}
The triple $(p;i,j,k)$ is \emph{admissible} if $p$ is prime, $p\ne3$,
$0<i<j<k$, $2k<p$, and $p\mid i^2+j^2+k^2$.
\end{definition}

\begin{theorem}[Admissible triple classification]\label{thm:main}
Suppose $(p;i,j,k)$ is admissible and $0\le R<p$.
Then $a_{p,v}(pN+R)=0$ for every $N\ge0$ if and only if there are
$w\in\Z^3$, $\lambda\in\Z$ and $e\in\{1,2\}$ such that
\[
 w\equiv\lambda v\pmod p,\qquad \norm{w}^2=ep,
\]
and, on setting $u=v\cdot w/p$ and $c=2/e$, one has
\begin{align}
 2\lambda R&\equiv w_1+w_2+w_3-6u\pmod p,
       &&ep\equiv1\pmod3,\label{eq:target1}\\
 c\lambda(2\lambda R-w_1-w_2-w_3+6u)&\equiv12\pmod p,
       &&ep\equiv2\pmod3.\label{eq:target2}
\end{align}
\end{theorem}
Here $u$ is an integer by isotropy, and $\lambda$ is automatically nonzero
modulo $p$. Indeed, $w\in p\Z^3$ and $\norm{w}^2=ep$ would imply
$p\mid e$, which is impossible. The residue equation is part of the
criterion: existence of a short root alone does not determine $R$.

The exact formal endpoint is
\lean{admissible_root_vanishing_equivalence} in
\lean{RamanujanTau/MultiQuintupleFullConverse.lean}. Its hypotheses are
precisely admissibility and $R<p$. Its vanishing predicate refers to the
coefficients of the actual infinite products, rather than a truncated
series or a hypothesized analytic identity.

The prime and distinct-index assumptions are material. For example,
common scaling at composite modulus can produce automatic support
vanishings with no root of the stated kind. The repository proves such
an obstruction at modulus 9 with three indices equal to 3. No claim about
all triples, composite moduli, or repeated indices is made here.

\section{Relation to the literature}
Daniels, Huber, McLaughlin and Ye study two-factor products by dissection
and lattice involutions \cite{dissection,vanishing}. Their concluding
sections discuss higher-factor examples; in particular, the second paper
lists the triple $(31;1,5,6)$ and the quadruple $(41;1,4,5,9)$ as
experimental examples, and announces further work. The theorem above
addresses an entire admissible class and includes necessity.
Their earlier discussion also records an already proved triple example
at modulus 13; individual triple identities are therefore not claimed
as a new phenomenon.

Chern and Tang give general coefficient-vanishing methods using theta
series and huffing operators \cite{chern-tang}. The distinction requiring
specialist review is whether the exact necessary-and-sufficient criterion
of Theorem~\ref{thm:main}, under its precise admissibility conditions,
already follows from that work or subsequent literature. Searches and
comparisons recorded on November 11, 2026 did not identify that exact
classification. This is evidence for further review, not a priority
certification. The present draft makes no ``first proof'' claim.

\section{Proof of the classification}
We give the argument with the intermediate arithmetic exposed. The
linked formal modules provide the complete details of convergence,
summation, lattice generation and the exceptional finite range.

\subsection{From coefficients to dual shells}
Watson's identity writes one factor as
\[
 Q(q^i,q^p)=\sum_{n\in\Z}\bigl(q^{A_i(n)}-q^{B_i(n)}\bigr),\quad
 A_i(n)=\frac{pn(3n-1)}2+3in,\quad B_i(n)=A_i(n)+i+pn.
\]
All these exponents are nonnegative for canonical indices. For a branch
$b\in\{0,1\}^3$, the progression condition becomes
$3v\cdot n+v\cdot b\equiv R\pmod p$.
Define
\[
 L=\{x\in\Z^3:v\cdot x\equiv0\pmod p\},\quad
 \mathcal N=pL^*=\{w\in\Z^3:w\equiv\lambda v\pmod p\}.
\]
Their covolumes are $p$ and $p^2$. Choose $q_0$ with
$3kq_0\equiv1\pmod p$, and set
\[
 c_b\equiv q_0(R-v\cdot b)\pmod p,\qquad
 x_b=\frac vp+\frac{2b-\mathbf1}{6}+c_b(0,0,1).
\]
Completing squares and applying Poisson summation to the eight Gaussian
coset sums shows that persistent vanishing forces the amplitude sum on
each norm shell of $\mathcal N$ to be zero. All sums involved are
absolutely convergent at positive Gaussian parameter.

For completeness, let $\zeta=e^{2\pi i/(6p)}$ and $u=v\cdot w/p$.
The dual amplitude is
\[
 A_R(w)=\zeta^{C(w)}\prod_{s=1}^3(1-\zeta^{\delta_s(w)}),\quad
 C(w)=6w_3q_0R+6u-\sum_s w_s,\quad
 \delta_s(w)=2w_s-6w_3q_0v_s.
\]
It is nonzero exactly when every coordinate of $w$ is nonzero modulo 3;
call such a vector supported. Write $t_s=\delta_s(w)/(2p)\pmod3$ and
let $n(w)$ be the number of $t_s$ equal to 2. For supported $w$,
\[
 A_R(w)=(1-\zeta^{2p})^3\rho^{a(w)},\qquad
 \rho=\zeta^3,\qquad a(w)=\frac{C(w)+pn(w)}3\pmod{2p}.
\]
The numerator is divisible by 3, and $a(-w)=p-a(w)$.
The cyclotomic relation for a primitive $2p$th root, combined with this
antipodal identity, forces equal multiplicities at phases $a$ and $a+p$.
Consequently every supported vector has an equal-norm opposite-phase
partner. This argument does not impose a bound on the size of a shell.

\subsection{A coherent pairing below the short-vector threshold}
Work modulo $p$, and define
\[
 h=\frac{v\cdot v}{p},\quad c_0=\frac{2R-\sum v_s}{3},\quad
 D=h+c_0,\quad
 \ell(w)=2\frac{v\cdot w}{p}+c_0\lambda(w).
\]
The scalar $\lambda(w)$ is uniquely determined modulo $p$.
Expanding $w=\lambda v+pz$ gives
\begin{align}
 \norm{w}^2/p&=\lambda(w)\ell(w)-D\lambda(w)^2,\label{eq:norm}\\
 2w\cdot z/p&=\lambda(w)\ell(z)+\lambda(z)\ell(w)
                   -2D\lambda(w)\lambda(z).\label{eq:pairing}
\end{align}
All pairings in $\mathcal N$ are divisible by $p$.

\begin{lemma}[Short-vector uniqueness]
If distinct integral vectors $w,w'$ have the same norm and are congruent
modulo an odd integer $p$, then $\norm{w}^2\ge p^2/2$.
\end{lemma}
\begin{proof}
Write $w'-w=pz$. Equality of norms gives
$2w\cdot z+p\norm{z}^2=0$, so $\norm{z}^2$ is a positive even integer.
The parallelogram identity gives $p^2\norm{z}^2\le4\norm{w}^2$.
\end{proof}

Let $S$ consist of supported vectors of norm squared less than $p^2/2$.
For opposite partners $w,w'\in S$, the phase condition preserves
$\ell$ modulo $p$. Their scalars are different by the lemma. Subtracting
\eqref{eq:norm} gives
\[
 D(\lambda(w)+\lambda(w'))=\ell(w).
\]
If $S$ is nonempty then $D\ne0$: otherwise $\ell(w)=0$ and
$\norm{w}^2/p\equiv0\pmod p$, contradicting
$0<\norm{w}^2<p^2/2$.
Thus partners are unique. Equation~\eqref{eq:pairing} shows that the
partner permutation preserves dot products modulo $p^2$.
Cauchy--Schwarz bounds the absolute difference of the two dot products
strictly below $p^2$, so it preserves them as integers.

\subsection{Integral generation, uniformly in the prime}
For $p\ne13$, $S$ generates $\mathcal N$ over $\Z$.
Here is the quantitative part of that assertion.
The integral lattice with underlying module $\mathcal N$ and inner
product $(w\cdot z)/p$ has Gram determinant $p$. Its adjugate lattice
is identified with $L$, whose minimum squared norm is at least 3:
norm 1 would force a zero index modulo $p$, and norm 2 would force two
indices equal up to sign.

A reduced basis has Gram--Schmidt squared lengths $A,B,C$ and
coefficients $\mu,\nu,\xi$ satisfying
\[
 |\mu|,|\nu|,|\xi|\le\tfrac12,\quad ABC=p,\quad
 B\ge\tfrac34A,\quad C\ge\tfrac34B,\quad
 A\in\Z_{>0},\quad M=AB\in\Z_{\ge3}.
\]
Take coefficient vectors $(x,y,z)$ with $x,y\in\{-2,-1,0,1,2\}$,
$z\in\{0,1\}$, which are nonzero isotropic vectors for the Gram form
modulo 3. Their images are supported and generate $\mathcal N$.
Indeed, reduction to $\F_3^3$ is an isomorphism, a nonzero isotropic
vector for a sum of three squares has all coordinates nonzero, and
the indicated box generates integrally: its reductions span, and
differences of representatives produce $3\Z^3$.
Their scaled squared norms are bounded by
\[
 C+\frac{49A+25B}{4}=\frac pM+\frac{49A+25B}{4}.
\]
If $M<12$, then $A\le3$, $B\le11$, giving $p/3+211/2<p/2$
for $p\ge1250$. If $M\ge12$, the reduction inequalities give
$A^3\le64p/27$ and $B^2\le4p/3$; hence $A\le p/80$ and
$B\le p/30$ for $p\ge1250$. The bound is then
$427p/960<p/2$. Returning to the unscaled norm gives the required
strict bound $p^2/2$.

Below 1250 there are 14,370 normalized canonical triples at 202 primes.
For 14,369 of them the supplied certificate consists of at most five
supported short vectors whose $3\times3$ minors have gcd $p^2$.
That gcd proves integral generation, rather than rational spanning.
Normalization, signed coordinate permutations and projective rescaling
cover every admissible triple in this range. The exception is modulus 13.
There the two canonical triples are $(1,3,4)$ and $(2,5,6)$.
Their zero residues are respectively $\{2,4,10\}$ and $\{3,9,11\}$;
six root certificates and twenty explicit nonzero coefficients settle
both directions. All finite certificates and the coverage argument are
checked in Lean. An independent integer-only checker is also supplied.

\subsection{Extracting the root and its residue}
For $p\ne13$, the dot-product-preserving permutation of the integral
generating set $S$ extends uniquely to a rational orthogonal involution
$T$, with $T(\mathcal N)=\mathcal N$ and hence $T(L)=L$.
It is neither $I$ nor $-I$. The first is excluded by the opposite-phase
condition. The second would force $\ell$ to vanish on $\mathcal N$,
whereas $\ell(pe_s)=2v_s\ne0\pmod p$.
In dimension three, $T$ is a reflection $H_w$ or its negative, for a
primitive integral normal $w$.

Set $n=w\cdot w$. Integrality on $pe_s\in L$ gives
$n\mid2p w_s w_t$ for every $s,t$, and primitivity implies $n\mid2p$.
The possibilities $n=1,2$ contradict admissibility. Thus $n=ep$,
$e\in\{1,2\}$; integrality on $L$ then gives $w\equiv\lambda v\pmod p$.

The phase parity on a supported vector $x$ is
\[
 (-1)^{a(x)}=(-1)^{\sum x_s}\chi_3(p)\chi_3(x_1x_2x_3).
\]
Every orthogonal matrix for a sum of three squares over $\F_3$ is a
signed permutation. Applying this to $T$ selects $T=H_w$ when
$ep\equiv1\pmod3$ and $T=-H_w$ when $ep\equiv2\pmod3$.
With $c=2/e$ and $t_x=x\cdot w/p$, the reflection is
$H_w(x)=x-ct_xw$. Preservation of $\ell$ gives $\ell(w)=0$ in
the first case and $c\lambda\ell(w)=4$ in the second.
Since $3\ell(w)=6u+\lambda(2R-\sum v_s)$, these are exactly
\eqref{eq:target1} and \eqref{eq:target2}.

Conversely, a root with the indicated residue equation induces the
sign-reversing involution on the Watson branches, preserving the
exponent in every coefficient sum. Pairwise cancellation proves
persistent vanishing. This direction is the formal theorem
\lean{hasProjectiveRootTarget_implies_persistentTripleVanishing}.
Together with the exceptional computation at 13, this proves
Theorem~\ref{thm:main}.

\section{An exact finite decision procedure}
\begin{proposition}\label{prop:scan}
For an admissible triple, all vanishing residues can be found by scanning
the $p-1$ nonzero scalars modulo $p$. For each scalar $\lambda$, take the
centered representative $w_s\in[-(p-1)/2,(p-1)/2]$ of $\lambda v_s$.
Retain $w$ if its squared norm is $p$ or $2p$, and solve the applicable
residue equation for $R$. The resulting set is exactly the set of
persistent zero progressions.
\end{proposition}
\begin{proof}
Admissibility implies $p\ge7$. If $\norm{w}^2\le2p$, every integral
coordinate satisfies $|w_s|\le(p-1)/2$, since
$((p+1)/2)^2>2p$ for $p\ge7$. Thus every root in
Theorem~\ref{thm:main} is the centered representative for its scalar.
Conversely, every retained vector satisfies the root conditions.
The coefficients $2\lambda$ and $2c\lambda^2$ in the two equations are
units modulo $p$, so each root gives one residue. Apply the theorem.
\end{proof}
This takes $O(p)$ arithmetic operations when inverses of all nonzero
residues are precomputed, and $O(p\log p)$ using a separate extended
Euclidean inversion for each retained root. This counts arithmetic
operations, not bit operations. No coefficient truncation is needed to
decide persistent vanishing. The supplied Python implementation uses
exact integers and modular inverses; it is a reference implementation,
not extracted Lean code. The completeness ingredient that every certified
root is the centered representative for its scalar is also checked in
Lean as \lean{projectiveRootTarget_centered}. The executable's full
input/output correspondence is the elementary argument above, rather than
a code-extraction or refinement theorem.

\section{A four-factor benchmark and the remaining converse}
For
\[
 G(q)=Q(q^5,q^{97})Q(q^4,q^{97})Q(q^9,q^{97})Q(q^{13},q^{97})
      =\sum_{n\ge0}g(n)q^n,
\]
the four known infinite vanishings arise from the reflection family
$(m-n,n,m,m+n)$ with $(m,n)=(9,4)$.
\begin{samepage}
\begin{theorem}[Complete classification at 97]\label{thm:97}
For $0\le R<97$,
\[
 (\forall N\ge0,\ g(97N+R)=0)
 \quad\Longleftrightarrow\quad R\in\{9,13,55,64\}.
\]
Every other residue has a nonzero coefficient with $N\le15$.
Moreover $g(97N+54)=0$ for $0\le N<15$, and $g(1509)=1$.
Consequently the common cutoff 15 is optimal for this product.
\end{theorem}
\end{samepage}
\begin{proof}
The reflection theorem proves the four infinite zero progressions.
For each of the other 93 residues, a generated exact coefficient
certificate proves nonvanishing by $N=15$. All required degrees are at
most 1509. Nonnegativity of the four Watson exponents and their quadratic
growth bound every contributing integer index in $[-3,3]$.
The resulting finite coefficient sums, the 15 preceding zero coefficients
in residue 54, and the nonzero value at 1509 are evaluated by Lean's kernel.
The Python generator only proposes certificates.
\end{proof}

The universal four-factor problem remains open here. In particular,
classifying compatibility within a collection of sixteen reflection
maps does not prove that every vanishing progression comes from one of
those maps. A successful general converse must establish that missing
necessity statement. A uniform quarter-prime first-witness bound cannot
be used for four factors: the repository proves counterexamples, including
one inside the reflection family at modulus 2521. The triple
quarter-prime conjecture is a separate unresolved quantitative question.

\section{Formal verification and review guide}
The reference snapshot for the existing triple proof is repository commit
\lean{8eff6d92d440bc3f6753abf32ca2a8afe187213b}, with Lean
\texttt{4.30.0-rc2} and the dependencies pinned by \lean{lake-manifest.json}.
The new four-factor theorem and this manuscript are tied to the source
hashes in \lean{docs/verification/quintuple_review.json}.
All public declarations in the dedicated audit must depend only on
\lean{propext}, \lean{Classical.choice} and \lean{Quot.sound}.
An axiom audit is accompanied by statement and definition checks;
axiom cleanliness by itself does not establish that a theorem states
the intended mathematical claim.

The principal source map, with paths relative to \lean{RamanujanTau/}, is:
\begin{center}\small
\begin{tabular}{>{\raggedright\arraybackslash}p{0.30\textwidth}>{\raggedright\arraybackslash}p{0.61\textwidth}}
\toprule
Mathematical step & Formal module (prefix \lean{MultiQuintuple}) \\
\midrule
Actual products and root predicate & \src{RootVanishingClassification.lean} \\
Poisson bridge and shell cancellation & \src{DualGaussianCancellation.lean} \\
Phase normalization and antipodes & \src{CanonicalPhaseCharacter.lean} \\
Opposite supported partners & \src{DualShellPartners.lean} \\
Integral short generation & \src{UniversalShortGeneration.lean} \\
Exceptional modulus 13 & \src{ConverseThirteen.lean} \\
Arithmetic extraction of the root & \src{ShortDualConverse.lean} \\
Final equivalence & \src{FullConverse.lean} \\
Centered-root completeness & \src{CenteredRoots.lean} \\
Four-factor converse at 97 & \src{FourNinetySevenConverse.lean} \\
\bottomrule
\end{tabular}
\end{center}
Run the reproduction script from the repository root:
\begin{verbatim}
python3 scripts/verify_quintuple_review.py
python3 scripts/classify_sparse_triple.py 31 1 5 6
\end{verbatim}
The first command builds the relevant modules, checks the dedicated axiom
audit, independently replays the finite lattice certificates, compares
the reference root scan with direct products on a small exhaustive census,
and checks the generated modulus-97 source and coefficient witnesses.
It writes a verification record and logs. Building the entire repository
is a separate broader check.

Independent review should concentrate on the exact product convention,
the normalization of the Fourier phase, passage from shell cancellation
to opposite partners, integral generation rather than real spanning,
the exhaustive normalized coverage below 1250, and extraction of the
two residue equations. The literature comparison must also account for
the further work announced in \cite{vanishing}.

\subsection*{Assistance and publication status}
This draft and the formal development were prepared with AI assistance.
The draft has not undergone independent human mathematical review or
peer review. It is prepared for that review; no journal submission,
external correspondence or priority claim is implied.

\begin{thebibliography}{9}
\bibitem{dissection}
T. Daniels, T. Huber, J. McLaughlin and D. Ye,
\emph{The $p$-Dissection of a Product of Quintuple Products},
arXiv:2603.04666 (2026),
\url{https://arxiv.org/abs/2603.04666}.
\bibitem{vanishing}
T. Daniels, T. Huber, J. McLaughlin and D. Ye,
\emph{Vanishing Coefficients in Products of Quintuple Products},
arXiv:2606.06701 (2026),
\url{https://arxiv.org/abs/2606.06701}.
\bibitem{chern-tang}
S. Chern and D. Tang,
\emph{General coefficient-vanishing results associated with theta series},
Advances in Applied Mathematics \textbf{159} (2024), 102742,
\url{https://arxiv.org/abs/2208.09882}.
\bibitem{repo}
T. DiFiore, \emph{RamanujanLean}, source and verification records,
\url{https://github.com/tomdif/RamanujanLean}.
\end{thebibliography}
\end{document}
